\documentclass[a4paper]{scrartcl}
\usepackage[
    typ=ueb,
    klausurtyp=kurs,
    klasse=scrartcl,
    fach=LF05 Physik,
    lerngruppe=T19E,
    loesungen=seite,
    module={Symbole,Bewertung}    
    %erwartungshorizontAnzeigen,
    %erwartungshorizontStil=simpel,
hHh]{schule}
% Dieses Dokument gehört zu den Beispiel des LaTeX Paketes Schule und ist von den Autoren
% des Pakets erstellt worden.
%
% Das Dokument steht unter der Lizenz: Creative Commons by-nc-sa Version 4.0
% http://creativecommons.org/licenses/by-nc-sa/4.0/deed.de
%
% Nach dieser Lizenz darf das Dokument beliebig kopiert und bearbeitet werden,
% sofern das Folgeprodukt wiederum unter gleichen Lizenzbedingungen vertrieben
% und auf die ursprünglichen Urheber verwiesen wird.
% Eine kommerzielle Nutzung ist ausdrücklich ausgeschlossen.

\author{Heiko Schröter}
\date{\today}
\title{Mechanik\\
(Arbeit, Leistung, Impuls und Stoß)}
\subtitle{Übungsaufgaben zur Arbeit, Leistung, Impuls und Stoß}

\usepackage{amsmath}
\usepackage{pgfplots}
\usepackage{mathrsfs}
\usepackage{xfrac}
\usepackage{siunitx}
%\sisetup{quotient-mode = fraction, locale = DE}
\sisetup{per-mode = fraction, locale = DE}
\usepackage{blindtext}
\usepackage{enumitem}
\newcommand{\Ergebnis}[1]{\underline{\underline{#1}}}

\begin{document}
%\punktuebersicht*

\setzeAufgabentemplate{schule-keinepunkte} 

\begin{aufgabe}[points={6}]
	Ein PKW mit der Masse \SI{1400}{\kilogram} und der Geschwindigkeit \SI{36}{\kilo\meter\per\hour} erhält von einem von hinten auffahrenden Wagen einen Kraftstoß \SI{2400}{\newton\second}.
	\begin{enumerate}[label=(\alph*)]
	\item Wie groß ist seine Geschwindigkeit unmittelbar nach dem Unfall?
	\item Welchen Kraftstoß erfuhr ein Insasse mit der Masse \SI{70}{\kilogram}?
	\item Welche Kraft wirkte auf das Fahrzeug, wenn es \SI{0,5}{\second} lang beschleunigt worden ist?
	\end{enumerate}
	\begin{tikzpicture}
		\draw[step=0.5cm,gray,very thin] (0,0) grid (15,8);
	\end{tikzpicture}
	
    \begin{loesung}
    \begin{flalign*}
    &v_v=\SI{36}{\kilo\meter\per\hour}=\SI{10}{\meter\per\second}&
    \end{flalign*}
    \begin{enumerate}[label=(\alph*)]
	\item
		\begin{flalign*}
		&F\cdot \Delta t=m\cdot v_n-m\cdot v_v=m\cdot(v_n-v_v)\Rightarrow\\
		&v_n-v_v=F\cdot\dfrac{\Delta t}{m}=\dfrac{\SI{2400}{\newton\second}}{\SI{1400}{\kilogram}}=\SI{1,71}{\meter\per\second}\Rightarrow\\
		&v_n=v_v+\SI{1,71}{\meter\per\second}=\SI{10}{\meter\per\second}+\SI{1,71}{\meter\per\second}=\SI{11,71}{\meter\per\second}=\Ergebnis{\SI{42,2}{\kilo\meter\per\hour}}&
		\end{flalign*}
	\item    
		\begin{flalign*}
		&F\cdot\Delta t=m\cdot(v_2-v_v)=\SI{70}{\kilogram}\cdot\SI{1,71}{\meter\per\second}=\Ergebnis{\SI{119,7}{\newton\second}}&
		\end{flalign*}
	\item
		\begin{flalign*}
		&F=a\cdot m=\dfrac{v_n-v_v}{\Delta t}\cdot m=\dfrac{\SI{1,71}{\meter\per\second}}{\SI{0,5}{\second}}\cdot \SI{1400}{\kilogram}=\Ergebnis{\SI{4800}{\newton}}\\
		&\text{oder}\quad F=\dfrac{F\cdot\Delta t}{\Delta t}=\dfrac{\SI{2400}{\newton\second}}{\SI{0,5}{\second}}=\Ergebnis{\SI{4800}{\newton}}&
		\end{flalign*}	
	\end{enumerate}
    \end{loesung}
\end{aufgabe}
\vspace{0.3cm}

\begin{aufgabe}[points={6}]
	Ein Eisenbahnwagen mit der Masse \SI{10}{\tonne} rollt mit einer Geschwindigkeit \SI{1,6}{\meter\per\second} gegen einen Prellbock, von dem er nach \SI{0,4}{\second} gleich schnell zurückprallt. Wie groß sind
	\begin{enumerate}[label=(\alph*)]
	\item die Impulsänderung,
	\item der Kraftstoß,
	\item die mittlere Kraft auf den Wagen bzw. auf den Prellbock?
	\end{enumerate}
	\begin{tikzpicture}
		\draw[step=0.5cm,gray,very thin] (0,0) grid (15,7);
	\end{tikzpicture}
	
    \begin{loesung}
        \begin{flalign*}
    &v_v=\SI{1,6}{\meter\per\second};\quad v_n=-\SI{1,6}{\meter\per\second};\quad v_n-v_v=-\SI{3,2}{\meter\per\second}&
    	\end{flalign*}
    \begin{enumerate}[label=(\alph*)]
    \item    
		\begin{flalign*}
		&\Delta p=m\cdot(v_n-v_v)=\SI{10000}{\kilogram}\cdot (-\SI{3,2}{\meter\per\second})=\Ergebnis{-\SI{32000}{\kilogram\meter\per\second}}&
		\end{flalign*}
	\item
		\begin{flalign*}
		&F\cdot\Delta t=\Delta p=-\SI{32000}{\kilogram\meter\per\second}=\Ergebnis{-\SI{32000}{\newton\second}}&
		\end{flalign*}
	\item
		\begin{flalign*}
		&\text{Wagen:}\quad F_W=\dfrac{F\cdot \Delta t}{\Delta t}=\dfrac{-\SI{32000}{\newton\second}}{\SI{0,4}{\second}}=\Ergebnis{-\SI{80000}{\newton}}\\
		&\text{Prellbock:}\quad F_P=\Ergebnis{+\SI{80000}{\newton}}&
		\end{flalign*}	
	\end{enumerate}
    \end{loesung}
\end{aufgabe}
\vspace{0.3cm}

\begin{aufgabe}[points={6}]
	Ein Hammer (Masse \SI{0,2}{\kilogram}) trifft mit der Geschwindigkeit \SI{8}{\meter\per\second} einen Nagel, der \SI{5}{\milli\second} lang in das Holz getrieben wird. Wie groß sind
	\begin{enumerate}[label=(\alph*)]
	\item die Impulsänderung des Hammers,
	\item der Kraftstoß auf den Hammer,
	\item die Kraft auf den Nagel?
	\end{enumerate}
	\begin{tikzpicture}
		\draw[step=0.5cm,gray,very thin] (0,0) grid (15,6);
	\end{tikzpicture}
	
    \begin{loesung}
    \begin{enumerate}[label=(\alph*)]
    \item    
		\begin{flalign*}
		&\Delta p=m\cdot v_n-m\cdot v_v=m\cdot (v_n-v_v)=\SI{0,2}{\kilogram}\cdot (0-\SI{8}{\meter\per\second})=\Ergebnis{-\SI{1,6}{\kilogram\meter\per\second}}&
		\end{flalign*}
	\item
		\begin{flalign*}
		&F\cdot\Delta t=\Delta p=-\SI{1,6}{\kilogram\meter\per\second}=\Ergebnis{-\SI{1,6}{\newton\second}}&
		\end{flalign*}
	\item
		\begin{flalign*}
		&\text{Nagel:}\quad F_N\cdot\Delta t=+\SI{1,6}{\newton\second}\Rightarrow \quad F_N=\dfrac{+\SI{1,6}{\newton\second}}{\SI{5}{\milli\second}}=\Ergebnis{\SI{320}{\newton}}&
		\end{flalign*}	
	\end{enumerate}
    \end{loesung}
\end{aufgabe}
\vspace{0.3cm}

\begin{aufgabe}[points={2}]
Ein Auto mit der Masse $m_1=\SI{800}{\kilogram}$ fährt frontal mit einer Geschwindigkeit $v=\SI{80}{\kilo\meter\per\hour}$ gegen einen Brückenpfeiler. Es verhält sich dabei vollkommen plastisch, was die totale Zerstörung des Autos zur folge hat. Der Brückenpfeiler hat eine Masse $m_2=\SI{15000}{\kilogram}$. Wie groß ist die hypothetische Geschwindigkeit $v$, mit der sich kurzzeitig beide Massen bewegen?\\
	\begin{tikzpicture}
		\draw[step=0.5cm,gray,very thin] (0,0) grid (15,6);
	\end{tikzpicture}
	
    \begin{loesung}    
		\begin{flalign*}
		v&=\dfrac{m_1\cdot v_1+m_2\cdot v_2}{m_1+m_2}=\dfrac{\SI{800}{\kilogram}\cdot \SI{80}{\kilo\meter\per\hour}+\SI{15000}{\kilogram}\cdot 0}{\SI{800}{\kilogram}+\SI{15000}{\kilogram}}\\
		&=\Ergebnis{\SI{4,051}{\kilo\meter\per\hour}}
		\end{flalign*}
    \end{loesung}
\end{aufgabe}
\vspace{0.3cm}

\begin{aufgabe}[points={2}]
Beim Schmieden eines Werkstückes wird ein \SI{1000}{\gram} Hammer verwendet. Die Hammergeschwindigkeit beim Auftreffen beträgt \SI{5}{\meter\per\second}, und der Schmiedeweg (Maß der Werkstückdeformation) beträgt \SI{0.75}{\milli\meter}. Dies bedeutet, dass der Hammer auf einem Weg von \SI{0.75}{\milli\meter} vollkommen abgebremst wird. Berechnen Sie die mittlere Hammerkraft.\\
	\begin{tikzpicture}
		\draw[step=0.5cm,gray,very thin] (0,0) grid (15,6);
	\end{tikzpicture}
	
    \begin{loesung}    
		\begin{flalign*}
		&F\cdot t=m\cdot \Delta v\\
		&\text{gleichmäßig beschleunigte Bewegung: }s=\dfrac{\Delta v\cdot t}{2}\rightarrow t=\dfrac{2s}{\Delta v}\\
		\rightarrow &F\cdot\dfrac{2s}{\Delta v}=m\cdot \Delta v\rightarrow F=\dfrac{m\cdot\Delta v^{2}}{2\cdot s}&\\
		&F=\dfrac{\SI{1}{\kilogram}\cdot(\SI{5}{\meter\per\second})^{2}}{2\cdot\SI{0.00075}{\meter}}=\Ergebnis{\SI{16667}{\newton}}&
		\end{flalign*}
    \end{loesung}
\end{aufgabe}
\vspace{0.3cm}

\begin{aufgabe}[points={6}]
Ein Rakete hat eine Masse von $m_2=\SI{200}{\tonne}$. Bei einer Ausströmgeschwindigkeit der Gase von $v=\SI{270}{\meter\per\second}$ wird in der Sekunde $\Delta m=\SI{300}{\kilogram}$ Raketentreibstoff verbrannt.
	\begin{enumerate}[label=(\alph*)]
	\item Stellen Sie eine Impulsbilanz für das System Rakete/Treibstoff auf, und zwar für den Zustand vor dem Start und für den Zustand nach dem Start, bei dem die Verbrennungsgase mit $v_1$ ausströmen.\\
	\begin{tabular}{|l|p{2.7cm}|p{2.7cm}|p{4.5cm}|}
    	    \hline 
    	    \rule[-1ex]{0pt}{2.5ex}   & Impuls Treibstoff & Impuls der Rakete & Gesamtimpuls des Systems Rakete/Treibstoff \\ 
    	    \hline 
    	    \rule[-1ex]{0pt}{2.5ex} vor dem Start & $p_1=$ & $p_2=$ & $p=$ \\ 
    	    \hline 
    	    \rule[-1ex]{0pt}{2.5ex} nach dem Start & $p_1=$ & $p_2=$ & $p=$ \\ 
    	    \hline 
    	    \end{tabular}
	\item Ermitteln Sie eine Funktion für die Geschwindigkeit $v_2=f(v_1,m_1,m_2)$, wobei zu Beginn Rakete und Treibstoff in Ruhe waren.
	\item Wie groß ist die Raketenbeschleunigung $a$?
	\end{enumerate}
	\begin{tikzpicture}
		\draw[step=0.5cm,gray,very thin] (0,0) grid (15,10);
	\end{tikzpicture}
	
    \begin{loesung}
    \begin{enumerate}[label=(\alph*)]
    \item
    	\begin{tabular}{|l|p{3cm}|p{3cm}|p{5cm}|}
    	    \hline 
    	    \rule[-1ex]{0pt}{2.5ex}   & Impuls Treibstoff & Impuls der Rakete & Gesamtimpuls des Systems Rakete/Treibstoff \\ 
    	    \hline 
    	    \rule[-1ex]{0pt}{2.5ex} vor dem Start & $p_1=m_1\cdot 0=0$ & $p_2=m_2\cdot 0=0$ & $p=p_1+p_2=0$ \\ 
    	    \hline 
    	    \rule[-1ex]{0pt}{2.5ex} nach dem Start & $p_1=m_1\cdot v_1$ & $p_2=m_2\cdot v_2$ & $p=p_1+p_2=m_1\cdot v_1+m_2\cdot v_2$ \\ 
    	    \hline 
    	    \end{tabular}     
	\item
		\begin{flalign*}
		&I=F\cdot t=\Delta p=m_1\cdot v_1+m_2\cdot v_2=0&
		\end{flalign*}
		(da in einem abgeschlossenen System der Gesamtimpuls konstant ist und dieser ja Null ist).\\
		Somit:
		\begin{flalign*}
		&m_1\cdot v_1=-m_2\cdot v_2\rightarrow v_2=-v_1\cdot\dfrac{m_1}{m_2}&
		\end{flalign*}
		Minuszeichen: $v_2$ ist $v_1$ entgegengerichtet
	\item
		\begin{flalign*}
		&v_2=-v_1\cdot\dfrac{m_1}{m_2}=-\SI{270}{\meter\per\second}\cdot\dfrac{\SI{0.3}{\tonne}}{\SI{200}{\tonne}}=-\SI{0.405}{\meter\per\second}&\\
		&a=\dfrac{\Delta v}{\Delta t}=\dfrac{-\SI{0.405}{\meter\per\second}}{\SI{1}{\second}}=\Ergebnis{-\SI{0.405}{\meter\per\square\second}}
		\end{flalign*}	
	\end{enumerate}
    \end{loesung}
\end{aufgabe}
\vspace{0.3cm}

\end{document}